\section{Detailed Experimental Configuration}
\label{app:exp-config}

This appendix reports the complete experimental configuration used in
Section~\ref{sec:Evaluation}. Unless explicitly varied by an experiment, the
data partition, model initialization, participating-client sequence, local
minibatches, malicious-client indices, and local-randomization seeds are shared
by all compared aggregation methods for a given dataset and seed. Method-specific
server step sizes are listed below and are frozen using only the validation set
before the corresponding attack experiments are executed.

\paragraph{Datasets and models.}
We evaluate CIFAR-10, FEMNIST, and Tiny-ImageNet. CIFAR-10 uses a CIFAR-style
ResNet-18 with a $3\times3$ stride-one stem, no initial max-pooling layer,
adaptive global average pooling, and a 10-class output layer (11,173,962
parameters). FEMNIST uses a 62-class CNN with two $3\times3$ convolutional
layers of 32 and 64 channels, one $2\times2$ max-pooling layer, a 256-unit
fully connected layer, and a 62-class output layer (3,246,270 parameters).
Tiny-ImageNet uses the same small-image ResNet-18 stem with a 200-class output
layer (11,271,432 parameters) and $64\times64$ inputs.

FEMNIST retains its native writer-disjoint partition. We construct a population
of 1,000 writer clients and sample 100 clients without replacement in every
communication round. CIFAR-10 and Tiny-ImageNet are partitioned across 100
population clients with a label-wise Dirichlet distribution with concentration
$\alpha=0.5$; 20 clients are sampled without replacement per round. The
partition seed is fixed to 1 across methods. We require at least 64 samples per
CIFAR-10 client and 32 samples per Tiny-ImageNet client.

\paragraph{Federated training.}
We use fixed local optimizer steps rather than local epochs because FEMNIST
writers and Dirichlet clients contain different numbers of examples. Hence,
no conversion from steps to an ambiguous fractional number of local epochs is
performed. The convergence configuration is summarized below.

\begin{table}[t]
  \centering
  \small
  \begin{tabular}{lccccp{4.8cm}}
    \toprule
    Dataset & $N$ & Batch & Local steps & Rounds & Optimizer and learning rate \\
    \midrule
    CIFAR-10 & 20 & 64 & 5 & 1,000 & Client SGD, learning rate 0.05,
    momentum 0.9, Nesterov enabled, weight decay $5\times10^{-4}$.
    Server step: 1.0 for FedAvg and $2\times10^{-4}$ for RAIN and SignSGD. \\
    FEMNIST & 100 & 32 & 1 & 1,000 & One minibatch gradient per selected
    writer; server SGD without momentum or weight decay. Server step: 0.1 for
    FedAvg, $2\times10^{-4}$ for RAIN, and $10^{-4}$ for SignSGD. \\
    Tiny-ImageNet & 20 & 16 & 2 & 1,000 & Client SGD, learning rate 0.025,
    momentum 0.9, Nesterov enabled, weight decay $5\times10^{-4}$.
    Server step: 1.0 for FedAvg, $2\times10^{-4}$ for RAIN, and
    $5\times10^{-5}$ for SignSGD and FLOD. \\
    \bottomrule
  \end{tabular}
  \caption{Convergence-training hyperparameters. $N$ is the number of clients
  selected in each communication round, not the total client population.}
  \label{tab:app-training-config}
\end{table}

CIFAR-10 uses cosine decay to a server step of $10^{-5}$ for RAIN and SignSGD;
the FedAvg server step remains 1.0. FEMNIST uses cosine minima of $10^{-3}$,
$10^{-5}$, and $10^{-5}$ for FedAvg, RAIN, and SignSGD, respectively.
Tiny-ImageNet uses the constant server steps shown in
Table~\ref{tab:app-training-config}. Robustness, privacy--utility, and
reference-mismatch experiments run for 200 communication rounds and use two
local steps per selected client; method-specific server steps are selected on
the reserved validation set, frozen before attack evaluation, and archived in
the resolved run configuration. We evaluate validation accuracy every 10
rounds in convergence runs (every 5 rounds in 200-round runs), save a
checkpoint every 25 rounds, and use CUDA automatic mixed precision.

\paragraph{Shuffle-DP parameters.}
Unless otherwise specified, the local randomization level is
$\varepsilon_0=10$ and $\delta=10^{-5}$. Each client update is clipped to
$\ell_2$ norm $C=1$ before Gaussian randomization. We use the explicit mapping
$\sigma/C=2/\varepsilon_0$, so the default noise multiplier is 0.2. The
end-to-end privacy budget is computed with the RDP accountant described in
Section~\ref{sec:Accounting}, using integer orders
$\Lambda=\{2,3,\ldots,128\}$. Experiments in
Fig.~\ref{fig:rain-privacy-tradeoff} vary $\varepsilon_0$ and recalibrate the
corresponding threshold; all other privacy parameters remain fixed.

\paragraph{Reference set and RAIN parameters.}
The audited reference set contains 124 FEMNIST examples, 500 CIFAR-10 examples,
and 400 Tiny-ImageNet examples. These examples are disjoint from client
training, threshold-calibration, validation, and test data. Threshold
calibration uses a further 3,200, 2,000, and 1,000 benign examples for FEMNIST,
CIFAR-10, and Tiny-ImageNet, respectively. For every dataset, partition, seed,
and $\varepsilon_0$, $\tau(\sigma)$ is fixed to the higher-interpolated 95th
percentile of normalized Hamming mismatch on this independent calibration
split and is never estimated from the current training round. Reference-size
and distribution-mismatch experiments vary only the stated reference-set
factor while holding the remaining split and training configuration fixed.

\paragraph{Attack configuration.}
The default robustness comparison uses 20\% Byzantine clients unless the
malicious-client fraction is the experimental variable. Krum Attack, Min-Max,
Scaling, and Attack-DPFL follow their original definitions. RAA and WOAA
construct the submitted sign messages as described in
Section~\ref{Attacks}. Within a dataset/seed/attack condition, all defenses
receive the same benign client updates before method-specific aggregation and
the same malicious-client indices and attacked submissions.

\paragraph{Baselines and hyperparameters.}
Baseline implementations follow the original papers or pinned author-released
code. Common data, optimization, clipping, randomization, and attack parameters
are shared across methods. FoundationFL uses its Median instantiation with one
synthetic update per selected client. DnC uses one filtering iteration with
$c=1$ and a 10,000-coordinate subsample. CONTRA uses selection fraction 1.0,
and FLARE uses at most one reference minibatch (32 FEMNIST, 64 CIFAR-10, or 16
Tiny-ImageNet examples). Defenses requiring a root/reference update receive the
same task-aligned audited reference split as RAIN. RFLPA's plaintext utility
experiment implements its cosine-trust aggregation rule; its packed-share
protocol cost is measured only with the pinned RFLPA artifact and is not
inferred from plaintext training time.

\paragraph{Hardware and implementation.}
Experiments run on one Ubuntu 22.04 LTS host with an Intel Core i9-14900KF
(24 CPUs exposed to the experiment process), 125~GiB usable RAM, and two
NVIDIA GeForce RTX 4090 D GPUs with 24,564~MiB memory each. The host has a
1~GbE interface. The artifact is a single-host protocol simulator: the two
logical non-colluding servers are isolated parties in one process, serialized
logical message bytes are counted exactly, and physical network latency and
TCP/IP framing are excluded from timed protocol measurements. GPU utility
training uses the plaintext realization of the stated aggregation function and
is not reported as MPC execution time.

Accuracy, balanced accuracy, loss, and attack-success results are reported as
mean and standard deviation over seeds 1, 2, and 3. Each protocol-efficiency
condition is measured five times after one untimed warm-up execution; we report
the mean and standard deviation and retain all five raw measurements. External
Camel and FLGuard measurements are executed on the same host in their pinned
environments, with their exact source commit and command stored alongside the
raw log.
